\section{What energy decay gives, and what it does not give}

Consider
\[
    u_t = \frac12 u u_{xx}, \qquad u\ge 0.
\]
Assume one has already proved that the energy
\[
    E(t):=\frac12\int_{\R}|u_x(t,x)|^2\,dx
\]
is decreasing and satisfies the identity
\[
    \frac{d}{dt}E(t)= -\frac12\int_{\R} u(t,x)|u_{xx}(t,x)|^2\,dx \le 0.
\]
Then the first conclusions are:
\begin{enumerate}[leftmargin=2em]
    \item $E(t)$ is monotone decreasing, hence has a limit $E_\infty\ge 0$;
    \item the dissipation is integrable in time,
    \[
        \int_0^{\infty}\int_{\R} u|u_{xx}|^2\,dx\,dt < \infty.
    \]
\end{enumerate}

This shows that the solution becomes asymptotically close to a steady regime in an averaged sense. However, this does \emph{not} by itself imply that $u(t,\cdot)$ converges to a single limit as $t\to\infty$.

\begin{proposition}
From energy decay alone, the natural next conclusion is a \emph{subsequence convergence} statement rather than a full convergence theorem.
\end{proposition}

Indeed, since the dissipation is integrable, there exists a sequence $t_n\to\infty$ such that
\[
    \int_{\R} u(t_n,x)|u_{xx}(t_n,x)|^2\,dx \to 0.
\]
Formally this means that along such times the solution is nearly stationary, because
\[
    u_t(t_n,\cdot)=\frac12 u(t_n,\cdot)u_{xx}(t_n,\cdot)
\]
is small in a weak sense.

If one can additionally prove compactness of the trajectory in a suitable topology, then one may extract a subsequence such that
\[
    u(t_n,\cdot) \to u_\infty
\]
for some limit profile $u_\infty$.

\begin{proposition}
Any such limit profile should satisfy the stationary equation
\[
    u_\infty (u_\infty)_{xx}=0
\]
in a weak sense.
\end{proposition}

Hence on every connected component of the positivity set $\{u_\infty>0\}$ one has
\[
    (u_\infty)_{xx}=0,
\]
so the limit profile is affine on each positive component. In a sufficiently regular globally bounded setting, this strongly suggests that any smooth global stationary limit must be constant.

\begin{remark}
What is still missing for full convergence is:
\begin{enumerate}[leftmargin=2em]
    \item precompactness of the whole trajectory $\{u(t,\cdot): t\ge 0\}$;
    \item uniqueness of possible stationary limit points.
\end{enumerate}
Without these two ingredients, energy decay only gives access to the $\omega$-limit set, not to convergence of the full orbit.
\end{remark}

On the whole line $\R$, this is exactly where difficulties appear. One does not automatically have the compactness mechanisms available on bounded domains, and it is possible in principle for mass or oscillation to drift toward spatial infinity. Therefore, from the assumptions ``nonnegative, bounded initial data'' together with energy monotonicity, one should not immediately claim uniform convergence on $\R$.

The clean conclusion is the following:
\begin{itemize}[leftmargin=2em]
    \item energy decay strongly suggests existence of subsequential limits;
    \item any subsequential limit should solve the stationary equation;
    \item to upgrade this to convergence to a unique limit, one needs extra compactness and a uniqueness principle for stationary states.
\end{itemize}
